\begin{textsample}{NEG Dim 1 – 2011\_c5 – Score 4.00 – 2011\_c5/enade\_2011\_tec\_constr\_f1a5d330d7.txt}  \label{ex:f1_neg_318}
CAPÍTULO 5 DISTRIBUIÇÃO DOS CONCEITOS 5.1 PANORAMA NACIONAL DA DISTRIBUIÇÃO DOS CONCEITOS A Tabela 5.1 apresenta a quantidade e distribuição de cursos de Tecnologia em Construção de Edifícios participantes do ENADE/2011, por faixa de conceito e Grande Região.
A diferença entre os cursos tabulados neste capítulo e no capítulo 2 são os cursos sem conceito, em princípio, aqueles sem alunos concluintes que participassem da prova.
Observando -se os dados da Tabela 5.1, nota -se que, dos 20 cursos participantes, sete ( 35,0 \% do total nacional ) classificaram -se com conceito 4, o modal.
Este foi também o conceito modal na região Nordeste ( quatro de sete cursos ), Sudeste ( dois dos seis cursos ) e Centro-Oeste ( o único curso participante ).
Na Área de Tecnologia em Construção de Edifícios um curso ficou sem conceito ( SC ).
O conceito com a segunda maior frequência foi 3, com cinco cursos dos 20 do país.
Em nível nacional, dois cursos obtiveram conceito 1, dois cursos foram avaliados com conceito 2 e três cursos alcançaram o conceito 5.
Na análise por região, a região Norte participou com dois cursos ou 10,0 \% do total nacional.
Nesta região, o conceito 1 foi atribuído a um curso e o outro foi avaliado com conceito 3.
A região Nordeste participou com sete cursos ou 35,0 \% do total nacional.
Como já comentado, destes, quatro cursos, 57,1 \% em termos regionais, obtiveram conceito 4, o conceito modal para a região.
Os conceitos 1, 2 e 3 foram atribuídos a um curso \textbf{cada}, com percentual de 14,2 \% para \textbf{cada} curso participante.
Nesta região nenhum curso alcançou o conceito 5 nem ficou sem conceito.
Dos seis cursos participantes da região Sudeste, dois ( 33,3 \% ) obtiveram conceito 3 e os outros dois ( 33,3 \% ) obtiveram o conceito 4, ambos modais.
Dos outros dois cursos da região, um curso ( 16,6 \% ) recebeu o conceito 5 e o outro ficou sem conceito.
A região Sul contou com quatro cursos ( 20,0 \% do total nacional ) distribuídos da seguinte forma : um curso recebeu o conceito 2, outro curso ( 25,0 \% ) o conceito 3, e dois cursos foram avaliados com conceito 4, o modal regional.
Nenhum dos cursos da região Sul ficou sem conceito.
Na região Centro-Oeste o único curso da Área de Tecnologia em Construção de Edifícios recebeu conceito 4. 5.2 CONCEITOS POR CATEGORIA ADMINISTRATIVA E POR GRANDE REGIÃO A Tabela 5.2 apresenta a distribuição dos cursos participantes do ENADE/2011 de Tecnologia em Construção de Edifícios, por Categoria Administrativa, de acordo com os conceitos por eles alcançados, segundo as Grandes Regiões brasileiras.
Dos 20 cursos participantes, 13 ( 65,0 \% ) eram ministrados em instituições públicas e sete ( 35,0 \% ), em privadas.
De acordo com as informações da Tabela 5.2, em termos nacionais três dos 20 cursos considerados receberam o conceito 5, sendo dois cursos de IES Pública e um de IES Privada.
Dos 13 cursos participantes de IES públicas, o conceito 4 foi o valor modal, atribuído a seis cursos ( 46,2 \% ).
Entre os demais cursos da rede pública, um curso obteve o conceito 1 ( 7,7 \% da categoria ) ; um o conceito 2 ( 7,7 \% ) ; três cursos ( 23,1 \% ) o conceito 3 e dois cursos ( 15,4 \% ) o conceito 5.
Nesta combinação de Categoria Administrativa, Conceitos e Grandes Regiões brasileiras nenhum curso ficou sem conceito.
Na rede privada, o conceito modal foi 3, com dois cursos ( 28,6 \% ) dos sete da categoria.
Pertencia a esta categoria o único curso de Tecnologia em Construção de Edifícios que ficou sem conceito.
Os outros quatro cursos ficaram distribuídos nos conceitos 1, 2, 4 e 5.
Na análise por região, observa -se que, na região Norte um curso pertencia a uma instituição pública e o outro a uma instituição privada.
O participante da rede pública foi avaliado com conceito 3 e o curso privado com conceito 1.
Na região Nordeste, apenas um curso pertencia à rede privada, 14,3 \% dos sete cursos da região.
Este curso privado recebeu conceito 4.
As instituições públicas dessa região participaram com seis cursos ( 85,7 \% ), dos quais três obtiveram o conceito modal, 4.
Os três demais foram avaliados com os conceitos 1, 2 e 3.
Nesta combinação de região e categoria nenhum curso ficou sem conceito ou alcançou conceito 5.
Na região Sudeste os seis cursos se distribuíam igualmente nas duas Categorias Administrativas.
Na rede pública, dois cursos ( 66,7 \% ) receberam o conceito 4 e um obteve conceito 5 ( 33,3 \% ). dois cursos ( 66,7 \% ) da rede privada receberam conceito 3 e um curso ficou sem conceito.
As instituições públicas na região Sul participaram com dois cursos e as privadas com os outros dois.
Na rede pública os cursos foram avaliados com conceitos 3 e 5.
Um dos cursos da rede privada recebeu conceito 2 e o outro o conceito 5.
Nenhum curso ficou sem conceito.
A região Centro-Oeste, como já se sabe, participou com apenas um curso.
Este, oferecido por uma IES Pública recebeu o conceito 4. 5.3 CONCEITOS POR ORGANIZAÇÃO ACADÊMICA E POR GRANDE REGIÃO Na Tabela 5.3 encontra -se a distribuição dos conceitos atribuídos aos cursos participantes do ENADE/2011 na Área de Tecnologia em Construção de Edifícios, por Organização Acadêmica, segundo as Grandes Regiões brasileiras.
Dos 20 cursos participantes, 15 cursos ( 75,0 \% ) eram oferecidos em Universidades, nenhum em Centro Universitário, e os cinco demais ( 25,0 \% ) eram de Faculdades.
De acordo com os dados apresentados, três cursos foram avaliados com conceito 5 e destes, dois eram vinculados a Universidades.
Apenas um curso em Universidades ficou sem conceito e o conceito 4 foi o modal, com seis cursos.
Os demais cursos avaliados receberam os conceitos : 1 ( um curso ), 2 ( um curso ), 3 ( quatro cursos ) e conceito 5 ( dois cursos, como já mencionado ).
Nas Faculdades, nenhum dos cinco cursos avaliados ficou sem conceito e estes cursos se distribuíram do conceito do 1 ao 5, um curso em \textbf{cada} conceito.
Considerando -se separadamente as regiões brasileiras, verifica -se que, na região Norte, um curso pertencia a Universidade e foi alocado ao conceito 3.
O outro da região participante do ENADE/2011 era de Faculdade e recebeu o com conceito 1.
Na região Nordeste, dos sete cursos participantes, seis cursos era de Universidades.
Para estes o conceito modal foi 4, com três cursos.
Os outros três cursos receberam, \textbf{cada} um, os conceitos 1, 2 e 3.
O único curso pertencente a uma Faculdade nesta região recebeu conceito 1.
Na região Sudeste, as Universidades concentraram quatro dos seis cursos da região.
Entre os cursos em Universidades, dois cursos receberam o conceito 4, modal, um curso recebeu o conceito 3 e o outro ficou sem conceito.
As Faculdades foram representadas por dois cursos, que se distribuíram nos conceitos 3 e 5.
Dos quatro cursos da região Sul que foram avaliados, três eram de Universidades, e dois deles receberam o conceito modal foi 5.
O outro curso recebeu o conceito 3.
O único curso participante na região Sul oriundo de Faculdade foi avaliado com o conceito 2.
O único curso da Área de Tecnologia em Construção de Edifícios da região Centro-Oeste pertencia a uma Universidade e recebeu o conceito 4.

% matched lemmas: cada
\end{textsample}
